\section{Degree-two trees and phase transport}
\label{sec:degree-two}

Let a pruned tree core have only degree-two variables.  Contract each
variable \(j\), adjacent to rows \(u\) and \(v\), into an edge between
those rows.  The contracted graph is a tree on \(I\).

\begin{theorem}[Explicit transported obstruction]\label{thm:transport}
Choose a root row \(r\) and set \(\alpha_r=1\).  Orient every contracted
edge away from the root.  If variable \(j\) connects parent row \(u\)
to child row \(v\), define
\begin{equation}\label{eq:transport}
 \alpha_v
 =
 -\frac{\overline{c_{uj}}}{\overline{c_{vj}}}\alpha_u.
\end{equation}
Then every \(\alpha_i\) is nonzero,
\[
 C^*\alpha=0,
 \qquad
 \Ker C^*=\C\alpha,
\]
and the value is
\begin{equation}\label{eq:tree-closed}
 \Gamma(b,C)
 =
 \frac{\left|\sum_{i\in I}\overline{\alpha_i}b_i\right|}
      {\max_{i\in I}|\alpha_i|}.
\end{equation}
\end{theorem}

\begin{proof}
There is a unique root path to every row, so \eqref{eq:transport} is
consistent and nonzero.  For the column \(j\) joining \(u,v\), it gives
\[
 \overline{c_{uj}}\alpha_u+
 \overline{c_{vj}}\alpha_v=0.
\]
Thus \(C^*\alpha=0\).  The core has one complex defect by
\cref{cor:one-defect-topology}, so \(\alpha\) spans the kernel.
\Cref{thm:one-defect} gives \eqref{eq:tree-closed}.
\end{proof}

\begin{proposition}[Recovery of primal variables]\label{prop:recover}
Let \(y\) be any optimal residual supplied by
\cref{thm:residual-face} and set \(d=y-b\).  The solution of
\(C\theta=d\) can be recovered by repeatedly removing a nonroot
analysis leaf \(\ell\).  If \(j\) is its incident variable and \(u\) is
the parent row, set
\[
 \theta_j=\frac{d_\ell}{c_{\ell j}},
 \qquad
 d_u\leftarrow d_u-c_{uj}\theta_j,
\]
then delete \(\ell,j\).  The last root equation holds automatically.
\end{proposition}

\begin{proof}
Each step solves the leaf equation exactly and substitutes it into its
parent equation.  After all nonroot rows are removed, the only possible
remaining inconsistency is in the root row.  But
\(\langle\alpha,d\rangle=0\) and the transported nonzero coefficients
express the root equation as the unique linear dependence among the
rows, so it also vanishes.
\end{proof}

\begin{corollary}[Weighted incidence columns]\label{cor:incidence}
Suppose every contracted edge \(e=\{u,v\}\) has column
\[
 C_e=w_e(e_u-e_v),
 \qquad w_e\ne0.
\]
Then \(\alpha\equiv1\) and
\begin{equation}\label{eq:conservation}
 \Gamma(b,C)=\left|\sum_{i\in I}b_i\right|.
\end{equation}
\end{corollary}

Equation \eqref{eq:conservation} is the unique conservation law of a
tree incidence matrix.  For general complex edge coefficients, the
ordinary sum is replaced by the literal phase-and-magnitude transported
sum in \eqref{eq:tree-closed}.

\begin{remark}[No cycle holonomy on a tree]
The recursion \eqref{eq:transport} never compares two paths to the same
row.  A cycle would impose a multiplicative consistency condition on
the transported phases and magnitudes.  That genuinely new phenomenon
is not part of the forest theorem.
\end{remark}

